\documentclass[11pt,a4paper]{article}

\usepackage[utf8]{inputenc}
\usepackage[T1]{fontenc}
\usepackage[margin=2.5cm]{geometry}
\usepackage{listings}
\usepackage{longtable}
\usepackage{enumitem}
\usepackage{array}

\lstset{
  basicstyle=\ttfamily\small,
  breaklines=true,
  frame=single,
  showstringspaces=false,
  columns=fullflexible,
  tabsize=2,
  xleftmargin=0pt,
  xrightmargin=0pt
}

\setlist[itemize]{leftmargin=1.2em}
\setlist[enumerate]{leftmargin=1.4em}

\newcommand{\statut}[1]{\textbf{Statut : #1}}

\title{Tailor Platform Web \\ Base de reference PHP, HTML, CSS et Bootstrap}
\author{Notes personnelles}
\date{\today}

\begin{document}

\maketitle
\tableofcontents
\newpage

% ============================================================
\section{Introduction}
% ============================================================

Ce document est le pendant web de ma base de reference Dart/Flutter. Il suit
la meme logique : il grandit avec le projet, il distingue ce qui est
"deja appris", "a consolider", "nouveau" ou "avance", et il ne repart jamais
de zero lors d'une mise a jour.

Le support principal utilise pour la partie orientee objet de PHP est le
tutoriel \textit{Programmez en oriente objet en PHP} (Victor Thuillier,
OpenClassrooms). La syntaxe presentee ici est cependant celle de PHP moderne
(8.x), avec les types stricts, ce qui differe legerement du tutoriel original
plus ancien mais respecte les memes principes de base.

Beaucoup de notions de programmation ont deja ete vues cote Dart. Quand
c'est le cas, ce document fait le lien plutot que de tout ré-expliquer
depuis zero.

% ============================================================
\section{Etat du projet : Tailor Platform Web}
% ============================================================

\subsection{Objectif}

La version web reproduit exactement le meme perimetre que la V1 mobile
Flutter :

\begin{center}
\textit{Le tailleur vend des modeles et le client en achete.}
\end{center}

Meme choix qu'en mobile pour cette V1 :

\begin{itemize}
  \item pas de base de donnees ;
  \item pas de backend externe (PHP est lui-meme le backend, mais sans SGBD) ;
  \item pas d'authentification reelle ;
  \item pas de paiement reel ;
  \item pas de livraison reelle.
\end{itemize}

A la place d'une base de donnees, les donnees vivent dans la session PHP
(\verb|$_SESSION|). C'est l'equivalent web des \verb|List| en memoire cote
Flutter : les donnees disparaissent si on ferme completement le navigateur
(la session expire).

\subsection{Design suivi}

Le site reprend la palette et l'esprit visuel d'une maquette fournie
(couleurs terracotta, foret, sable, creme, typographie serif pour les
titres). L'implementation reste volontairement simple : HTML, CSS et
Bootstrap, sans framework JavaScript.

\subsection{Structure du projet}

\begin{lstlisting}
tailor_platform_web/
  src/
    Model/
      Tailor.php
      Realisation.php
    Support/
      Panier.php
      CartService.php
    Data/
      CatalogueRepository.php
      TailorRepository.php
  includes/
    bootstrap.php
    header.php
    footer.php
  public/
    index.php
    tailor_profile.php
    add_article.php
    catalogue.php
    article_detail.php
    cart.php
    order_confirmation.php
    assets/
      css/style.css
      vendor/bootstrap/css/bootstrap.min.css
      vendor/bootstrap/js/bootstrap.bundle.min.js
\end{lstlisting}

\verb|public/| est la racine web (document root). \verb|src/| et
\verb|includes/| contiennent du code PHP qui n'est jamais accede
directement par le navigateur : il est charge par les pages de
\verb|public/|.

\subsection{Parcours V1 (identique au mobile)}

\begin{lstlisting}
TAILLEUR
profil/catalogue
  |
ajouter un modele
  |
modele disponible dans le catalogue
  |
CLIENT
catalogue
  |
detail du modele
  |
ajouter au panier
  |
panier
  |
calcul du total
  |
commander
  |
confirmation de commande
  |
retour au catalogue
\end{lstlisting}

\subsection{Etat d'avancement}

L'ensemble du parcours ci-dessus est fonctionnel des la premiere version du
site web. Contrairement a la version Flutter (developpee etape par etape),
la version web a ete construite d'un bloc en reprenant les regles deja
etablies cote mobile.

% ============================================================
\section{PHP : vue d'ensemble des notions}
% ============================================================

\begin{longtable}{|p{5cm}|p{4cm}|p{5cm}|}
\hline
\textbf{Notion} & \textbf{Statut} & \textbf{Remarque} \\
\hline
Variables (\$), types & Nouveau & Similaire a Dart, mais prefixe \$ et typage plus souple \\
\hline
echo, affichage & Nouveau & \\
\hline
Tableaux indexes et associatifs & Nouveau & Jouent le role des List et Map de Dart \\
\hline
foreach & Nouveau & Equivalent du for-in de Dart \\
\hline
Conditions, operateurs & Deja appris & Meme logique qu'en Dart \\
\hline
Fonctions & Deja appris & Meme logique qu'en Dart, syntaxe differente \\
\hline
Classes, objets, \$this & Nouveau & Bases deja connues via Dart, syntaxe PHP a apprendre \\
\hline
Visibilite (private, public, protected) & Nouveau & Rapproche du principe d'encapsulation \\
\hline
Constructeur \_\_construct & Nouveau & Equivalent du constructeur Dart \\
\hline
Namespaces & Nouveau & Organise les classes, evite les conflits de noms \\
\hline
Autoload (spl\_autoload\_register) & Nouveau & Charge les classes automatiquement \\
\hline
Superglobales (\$\_GET, \$\_POST, \$\_SESSION, \$\_SERVER) & Nouveau & Coeur de la communication entre pages en PHP \\
\hline
Sessions PHP & Nouveau & Remplace la base de donnees pour cette V1 \\
\hline
include / require & Nouveau & Assemble les pages a partir de morceaux communs (header, footer) \\
\hline
header() et redirection & Nouveau & Equivalent web de Navigator.push/pop \\
\hline
Traitement de formulaire (POST) & Nouveau & \\
\hline
Validation de donnees utilisateur & Nouveau & \\
\hline
Exceptions (try/catch) & Avance & Pas encore utilise dans le projet \\
\hline
Interfaces, heritage, abstract & Avance & Vus dans le tutoriel PHP, pas encore utilises dans le projet \\
\hline
\end{longtable}

\subsection{Classes en PHP : Realisation et Tailor}

\statut{Nouveau}

\paragraph{C'est quoi.}
Comme en Dart, une classe decrit un modele de donnee et un objet en est une
instance. La grande difference de syntaxe : les attributs sont precedes de
\verb|$| et on accede a un attribut ou une methode avec \verb|->| au lieu de
\verb|.|.

\paragraph{Comparaison rapide avec Dart.}

\begin{longtable}{|p{4cm}|p{5cm}|p{5cm}|}
\hline
\textbf{Concept} & \textbf{Dart} & \textbf{PHP} \\
\hline
Definir une classe & \verb|class Realisation { }| & \verb|class Realisation { }| \\
\hline
Attribut & \verb|String title;| & \verb|private string $title;| \\
\hline
Constructeur & \verb|Realisation({required this.title})| & \verb|public function __construct(string $title)| \\
\hline
Acceder a un attribut/methode & \verb|objet.title| & \verb|$objet->getTitle()| \\
\hline
Soi-meme & \verb|this| & \verb|$this| \\
\hline
\end{longtable}

\paragraph{Syntaxe minimale.}

\begin{lstlisting}
<?php

declare(strict_types=1);

namespace App\Model;

class Realisation
{
    private int $id;
    private string $title;
    private string $description;
    private float $price;

    public function __construct(int $id, string $title, string $description, float $price)
    {
        $this->id = $id;
        $this->title = $title;
        $this->description = $description;
        $this->price = $price;
    }

    public function getId(): int
    {
        return $this->id;
    }

    public function getTitle(): string
    {
        return $this->title;
    }

    public function getDescription(): string
    {
        return $this->description;
    }

    public function getPrice(): float
    {
        return $this->price;
    }
}
\end{lstlisting}

\paragraph{Pourquoi des attributs private avec des getters.}
C'est le principe d'encapsulation expliquer dans le tutoriel PHP OOP :
personne en dehors de la classe ne doit pouvoir modifier directement
\verb|$title| ou \verb|$price|. On passe obligatoirement par une methode
publique (ici un getter) pour lire la valeur. Cela protege l'objet contre
des modifications incoherentes.

\paragraph{Creer un objet.}

\begin{lstlisting}
$realisation = new Realisation(0, 'Costume traditionnel', 'Costume deux pieces sur mesure', 45000.0);

echo $realisation->getTitle();
\end{lstlisting}

\paragraph{Erreurs frequentes.}
\begin{itemize}
  \item Oublier \verb|$this->| dans le constructeur : sans lui, la variable
  reste locale a la methode et l'attribut de l'objet n'est jamais rempli.
  \item Declarer un attribut \verb|public| par facilite, ce qui casse
  l'encapsulation.
  \item Oublier le point-virgule a la fin d'une instruction (erreur
  frequente venant de Dart ou l'oubli est parfois moins visible).
\end{itemize}

\paragraph{Dans Tailor Platform.}
\verb|App\Model\Realisation| et \verb|App\Model\Tailor| sont la copie
conforme, en PHP, des classes Dart du meme nom. Elles n'ont aucune logique
de session ou d'affichage : elles ne font que representer une donnee.

\subsection{Namespace et autoload}

\statut{Nouveau}

\paragraph{C'est quoi.}
Un \verb|namespace| est un peu comme un dossier logique pour les classes.
Il evite qu'une classe \verb|Realisation| du projet entre en conflit avec
une classe \verb|Realisation| d'une bibliotheque externe.

\paragraph{Syntaxe minimale.}

\begin{lstlisting}
namespace App\Model;

class Realisation
{
    // ...
}
\end{lstlisting}

Pour utiliser la classe ailleurs :

\begin{lstlisting}
use App\Model\Realisation;

$realisation = new Realisation(0, 'Titre', 'Description', 1000.0);
\end{lstlisting}

\paragraph{L'autoload.}
Sans autoload, il faudrait ecrire un \verb|require| pour chaque fichier de
classe avant de l'utiliser. L'autoload automatise ce chargement : des qu'une
classe inconnue est utilisee, PHP appelle une fonction qui va chercher le
bon fichier.

\begin{lstlisting}
spl_autoload_register(static function (string $class): void {
    $prefix = 'App\\';

    if (!str_starts_with($class, $prefix)) {
        return;
    }

    $relativeClass = substr($class, strlen($prefix));
    $file = __DIR__ . '/../src/' . str_replace('\\', '/', $relativeClass) . '.php';

    if (is_file($file)) {
        require_once $file;
    }
});
\end{lstlisting}

\paragraph{Point important rencontre dans le projet.}
L'autoload doit etre enregistre \textbf{avant} \verb|session_start()|. En
effet, PHP recharge automatiquement les objets stockes en session des le
demarrage de la session. Si la classe correspondante n'est pas encore
connue a ce moment (autoload pas encore enregistre), l'objet recupere est
incomplet et inutilisable. C'est une erreur reelle que nous avons rencontree
et corrigee dans \verb|includes/bootstrap.php|.

\paragraph{Dans Tailor Platform.}
Toutes les classes du projet vivent sous le namespace \verb|App\| (par
exemple \verb|App\Model\Realisation|, \verb|App\Data\CatalogueRepository|).
L'autoload est configure une seule fois, dans \verb|includes/bootstrap.php|,
et ce fichier est inclus en tout premier dans chaque page.

\subsection{Superglobales : \$\_GET, \$\_POST, \$\_SESSION, \$\_SERVER}

\statut{Nouveau}

\paragraph{C'est quoi.}
Les superglobales sont des tableaux automatiquement remplis par PHP, et
accessibles depuis n'importe quel fichier sans rien importer.

\begin{longtable}{|p{2.5cm}|p{9.5cm}|}
\hline
\textbf{Superglobale} & \textbf{Role} \\
\hline
\verb|$_GET| & Donnees envoyees dans l'URL (\verb|?id=3|) \\
\hline
\verb|$_POST| & Donnees envoyees par un formulaire en methode POST \\
\hline
\verb|$_SESSION| & Donnees qui persistent d'une page a l'autre pour un meme visiteur \\
\hline
\verb|$_SERVER| & Informations sur la requete en cours (methode HTTP, etc.) \\
\hline
\end{longtable}

\paragraph{Exemple : lire un identifiant dans l'URL.}

\begin{lstlisting}
$id = isset($_GET['id']) && is_numeric($_GET['id']) ? (int) $_GET['id'] : null;
\end{lstlisting}

\paragraph{Exemple : lire un champ de formulaire.}

\begin{lstlisting}
$title = trim((string) ($_POST['title'] ?? ''));
\end{lstlisting}

\paragraph{Exemple : connaitre la methode HTTP de la requete.}

\begin{lstlisting}
if ($_SERVER['REQUEST_METHOD'] === 'POST') {
    // le formulaire a ete soumis
}
\end{lstlisting}

\paragraph{Erreurs frequentes.}
\begin{itemize}
  \item Lire \verb|$_POST['title']| sans verifier qu'il existe : erreur si
  le champ n'a pas ete envoye. On utilise \verb|??| pour donner une valeur
  par defaut.
  \item Faire confiance aveuglement a une donnee venant de \verb|$_GET| ou
  \verb|$_POST| : il faut toujours la valider avant de l'utiliser.
  \item Oublier \verb|(int)| ou \verb|(float)| : ces donnees arrivent
  toujours sous forme de texte (\verb|string|), meme si elles ressemblent a
  des nombres.
\end{itemize}

\paragraph{Dans Tailor Platform.}
\verb|article_detail.php| lit l'identifiant du modele via \verb|$_GET['id']|.
\verb|add_article.php| lit le titre, la description et le prix via
\verb|$_POST|. \verb|cart.php| declenche la commande en verifiant
\verb|$_SERVER['REQUEST_METHOD']|.

\subsection{Sessions PHP}

\statut{Nouveau}

\paragraph{C'est quoi.}
Une session permet de conserver des donnees d'une page a l'autre pour un
meme visiteur, alors que chaque requete HTTP est normalement independante
des autres.

\paragraph{Comment la visualiser.}

\begin{lstlisting}
Visiteur A                      Visiteur B
$_SESSION (A)                   $_SESSION (B)
  cart: [Costume]                 cart: []
  realisations: [...]             realisations: [...]
\end{lstlisting}

Chaque visiteur a sa propre session, identifiee par un cookie invisible
envoye automatiquement par le navigateur.

\paragraph{Demarrer une session.}

\begin{lstlisting}
if (session_status() !== PHP_SESSION_ACTIVE) {
    session_start();
}
\end{lstlisting}

\paragraph{Lire et ecrire dans la session.}

\begin{lstlisting}
$_SESSION['cart'][] = $realisation;   // ajouter
$items = $_SESSION['cart'];           // lire
$_SESSION['cart'] = [];               // vider
\end{lstlisting}

\paragraph{Ce que je dois retenir.}
\verb|$_SESSION| se comporte comme un tableau normal, mais son contenu
survit d'une page a l'autre. C'est ce qui remplace la base de donnees dans
cette V1 : le "catalogue" et le "panier" sont juste des tableaux stockes en
session.

\paragraph{Dans Tailor Platform.}
\verb|CatalogueRepository| stocke la liste des \verb|Realisation| dans
\verb|$_SESSION['realisations']|. \verb|CartService| stocke le panier dans
\verb|$_SESSION['cart']|. Si le tailleur ajoute un modele, il est visible
tant que la session dure, exactement comme le \verb|cart.add()| en memoire
cote Flutter.

\subsection{header() et redirection}

\statut{Nouveau}

\paragraph{C'est quoi.}
\verb|header('Location: ...')| dit au navigateur : va sur cette autre page.
C'est l'equivalent web de \verb|Navigator.push| ou \verb|Navigator.pop| en
Flutter, mais cote serveur.

\paragraph{Syntaxe minimale.}

\begin{lstlisting}
header('Location: tailor_profile.php');
exit;
\end{lstlisting}

\verb|exit| est indispensable juste apres : sans lui, le reste du script
continue de s'executer meme si le navigateur va deja changer de page.

\paragraph{Passer une donnee a la page suivante.}
Cote Flutter, on passe une donnee via le constructeur de la page. Cote web,
l'equivalent le plus simple est de l'ajouter dans l'URL :

\begin{lstlisting}
header('Location: order_confirmation.php?total=' . rawurlencode((string) $total));
exit;
\end{lstlisting}

Et on la relit sur la page d'arrivee avec \verb|$_GET| :

\begin{lstlisting}
$total = isset($_GET['total']) ? (float) $_GET['total'] : 0.0;
\end{lstlisting}

\paragraph{Erreurs frequentes.}
\begin{itemize}
  \item Oublier \verb|exit| apres \verb|header('Location: ...')|.
  \item Appeler \verb|header()| apres avoir deja affiche du HTML : PHP
  refuse alors d'envoyer l'en-tete (erreur "headers already sent").
  \item Oublier \verb|rawurlencode| sur une valeur qui pourrait contenir des
  caracteres speciaux dans l'URL.
\end{itemize}

\paragraph{Dans Tailor Platform.}
\verb|add_article.php| redirige vers \verb|tailor_profile.php| apres un
ajout reussi. \verb|article_detail.php| redirige vers \verb|catalogue.php|
apres l'ajout au panier. \verb|cart.php| redirige vers
\verb|order_confirmation.php| en transmettant le total et le nombre
d'articles dans l'URL.

\subsection{Traitement d'un formulaire}

\statut{Nouveau}

\paragraph{Principe general.}
Une page qui contient un formulaire fait souvent deux choses a la fois :
elle affiche le formulaire (methode GET) et elle traite son envoi (methode
POST), dans le meme fichier PHP.

\paragraph{Structure type utilisee dans le projet.}

\begin{lstlisting}
<?php

if ($_SERVER['REQUEST_METHOD'] === 'POST') {
    // 1. recuperer les valeurs envoyees
    // 2. les valider
    // 3. si tout est correct, enregistrer puis rediriger
}

// 4. si on arrive ici, afficher le formulaire (premiere visite ou erreur)
?>

<form method="post">
  <input type="text" name="title">
  <button type="submit">Enregistrer</button>
</form>
\end{lstlisting}

\paragraph{Ce que je dois retenir.}
Le fait de re-afficher le formulaire (au lieu de rediriger) quand il y a une
erreur permet de montrer les messages d'erreur au visiteur, avec ce qu'il
avait deja tape.

\paragraph{Dans Tailor Platform.}
\verb|add_article.php| suit exactement cette structure : verification du
titre, de la description et du prix, affichage des erreurs si besoin, sinon
enregistrement via \verb|CatalogueRepository::add()| puis redirection.

% ============================================================
\section{HTML, CSS et Bootstrap : vue d'ensemble}
% ============================================================

\begin{longtable}{|p{5cm}|p{4cm}|p{5cm}|}
\hline
\textbf{Notion} & \textbf{Statut} & \textbf{Remarque} \\
\hline
Structure HTML de base & Nouveau & doctype, html, head, body \\
\hline
Formulaires HTML (form, input, textarea) & Nouveau & \\
\hline
CSS : selecteurs, proprietes & Nouveau & \\
\hline
CSS : variables (--ma-couleur) & Nouveau & Utilisees pour la palette du projet \\
\hline
Bootstrap : grille (container, row, col) & Nouveau & Equivalent visuel de Row/Column/Expanded en Flutter \\
\hline
Bootstrap : composants (card, navbar, btn, form-control) & Nouveau & \\
\hline
JavaScript & Avance & Pas encore utilise dans le projet (Bootstrap est inclus pour plus tard) \\
\hline
\end{longtable}

\subsection{Grille Bootstrap : container, row, col}

\statut{Nouveau}

\paragraph{C'est quoi.}
La grille Bootstrap permet d'organiser une page en lignes et colonnes, un
peu comme \verb|Row| et \verb|Column| en Flutter, mais pensee pour
s'adapter automatiquement a la largeur de l'ecran (responsive).

\paragraph{Comment la visualiser.}

\begin{lstlisting}
container
  row
    col   col   col
   [ A ] [ B ] [ C ]
\end{lstlisting}

Sur un ecran etroit (mobile), les colonnes peuvent se replacer les unes
sous les autres automatiquement.

\paragraph{Syntaxe minimale.}

\begin{lstlisting}
<div class="container">
  <div class="row">
    <div class="col">Colonne 1</div>
    <div class="col">Colonne 2</div>
  </div>
</div>
\end{lstlisting}

\paragraph{Classes utiles rencontrees dans le projet.}

\begin{longtable}{|p{4cm}|p{9cm}|}
\hline
\textbf{Classe} & \textbf{Role} \\
\hline
\verb|container| & Centre le contenu et lui donne une largeur maximale \\
\hline
\verb|row| & Cree une ligne de colonnes \\
\hline
\verb|col| & Une colonne de largeur egale et automatique \\
\hline
\verb|row-cols-1 row-cols-sm-2 row-cols-lg-3| & Nombre de colonnes different selon la taille d'ecran \\
\hline
\verb|g-4| & Espace (gap) entre les colonnes et les lignes \\
\hline
\verb|d-flex| & Active flexbox sur l'element \\
\hline
\verb|justify-content-between| & Espace les elements flexbox aux extremites \\
\hline
\verb|align-items-center| & Centre verticalement les elements flexbox \\
\hline
\end{longtable}

\paragraph{Erreurs frequentes.}
\begin{itemize}
  \item Mettre une classe \verb|row| directement sur le \verb|body|, sans
  \verb|container| autour.
  \item Mettre du contenu directement dans \verb|row| sans passer par
  \verb|col|.
  \item Oublier \verb|g-4| (ou une classe de gap) et se retrouver avec des
  cartes collees les unes aux autres.
\end{itemize}

\paragraph{Dans Tailor Platform.}
Le catalogue client et le catalogue du tailleur utilisent
\verb|row row-cols-1 row-cols-sm-2 row-cols-lg-3 g-4| pour afficher les
modeles en grille responsive, ce qui est l'equivalent visuel du
\verb|GridView.builder| cote Flutter.

\subsection{Composant Card}

\statut{Nouveau}

\paragraph{C'est quoi.}
Une \verb|card| Bootstrap est un bloc visuel avec un fond, une bordure et un
espacement interieur, utilise pour presenter une information de maniere
groupee (un modele, un produit, etc.).

\paragraph{Syntaxe minimale.}

\begin{lstlisting}
<div class="card">
  <div class="card-body">
    <h3 class="card-title">Titre</h3>
    <p class="card-text">Description</p>
  </div>
</div>
\end{lstlisting}

\paragraph{Dans Tailor Platform.}
Chaque modele du catalogue est affiche dans une \verb|card|, avec le titre,
la description, le prix et un bouton "Voir". C'est l'equivalent visuel du
\verb|Card| Flutter utilise dans le \verb|GridView.builder| mobile.

\subsection{Formulaires Bootstrap}

\statut{Nouveau}

\paragraph{Syntaxe minimale.}

\begin{lstlisting}
<div class="mb-3">
  <label for="title" class="form-label">Titre</label>
  <input type="text" class="form-control" id="title" name="title">
</div>
\end{lstlisting}

\paragraph{Classes utiles.}

\begin{longtable}{|p{3.5cm}|p{8.5cm}|}
\hline
\textbf{Classe} & \textbf{Role} \\
\hline
\verb|form-label| & Style du texte au-dessus du champ \\
\hline
\verb|form-control| & Style standard d'un champ de saisie \\
\hline
\verb|mb-3| & Marge en dessous, pour espacer les champs entre eux \\
\hline
\verb|btn btn-primary| & Bouton principal (couleur d'accent du site) \\
\hline
\verb|btn btn-outline-secondary| & Bouton secondaire, moins visible \\
\hline
\verb|alert alert-danger| & Bloc de message d'erreur \\
\hline
\end{longtable}

\paragraph{Dans Tailor Platform.}
Le formulaire \verb|add_article.php| utilise ces classes pour le titre, la
description et le prix, avec un bloc \verb|alert alert-danger| affiche en
cas d'erreur de validation.

% ============================================================
\section{Dictionnaire des fiches PHP}
% ============================================================

\subsection{Panier (classe)}

\statut{Nouveau}

\paragraph{C'est quoi.}
Une classe qui encapsule la logique du panier : ajouter un article,
compter les articles, calculer le total, vider le panier.

\paragraph{Pourquoi une classe plutot qu'un simple tableau.}
On pourrait se contenter d'un tableau de \verb|Realisation| en session,
comme cote Flutter (\verb|List<Realisation> cart = []|). Mais en
orientant objet, on regroupe les donnees (les articles) et les actions
possibles sur ces donnees (add, getTotal, clear) dans une seule classe. Cela
evite de recalculer le total a plusieurs endroits differents du code avec
des risques d'erreur.

\paragraph{Syntaxe minimale.}

\begin{lstlisting}
class Panier
{
    private array $items;

    public function __construct(array $items = [])
    {
        $this->items = $items;
    }

    public function add(Realisation $realisation): void
    {
        $this->items[] = $realisation;
    }

    public function getTotal(): float
    {
        $total = 0.0;

        foreach ($this->items as $realisation) {
            $total += $realisation->getPrice();
        }

        return $total;
    }
}
\end{lstlisting}

\paragraph{Comparaison avec Dart.}
La methode \verb|getTotal()| ressemble volontairement a la fonction
\verb|calculerTotal()| ecrite cote Flutter : meme logique (une boucle qui
accumule les prix), juste rangee a l'interieur d'une classe plutot que
laissee comme fonction independante.

\paragraph{Dans Tailor Platform.}
\verb|CartService| fait le lien entre cette classe \verb|Panier| et la
session PHP : il reconstruit un \verb|Panier| a partir de
\verb|$_SESSION['cart']| a chaque page, pour pouvoir utiliser ses methodes.

\subsection{Repository (CatalogueRepository, TailorRepository)}

\statut{Nouveau}

\paragraph{C'est quoi.}
Un repository est une classe dont le role est uniquement de fournir des
donnees (les recuperer, les ajouter), sans se soucier de l'affichage. C'est
une facon d'organiser le code pour que chaque classe ait une seule
responsabilite claire.

\paragraph{Dans Tailor Platform.}
\verb|CatalogueRepository| initialise le catalogue fictif si la session est
vide, retrouve un modele par son identifiant, et ajoute un nouveau modele.
\verb|TailorRepository| fournit le tailleur unique de cette V1. Ces deux
classes seront les premieres a changer le jour ou une vraie base de donnees
sera ajoutee au projet : le reste du code (les pages) n'aura pas besoin de
changer, puisqu'il ne fait qu'appeler leurs methodes.

% ============================================================
\section{Cheat sheet rapide}
% ============================================================

\subsection{Reference ultra rapide}

\begin{longtable}{|l|l|}
\hline
\textbf{Element} & \textbf{Role en un mot} \\
\hline
\$\_GET & lire l'URL \\
\hline
\$\_POST & lire un formulaire \\
\hline
\$\_SESSION & memoriser entre les pages \\
\hline
header('Location: ...') & rediriger \\
\hline
class & definir un modele de donnee \\
\hline
\$this & soi-meme, dans une classe \\
\hline
namespace & organiser les classes \\
\hline
container / row / col & grille Bootstrap \\
\hline
card & bloc visuel Bootstrap \\
\hline
form-control & champ de saisie stylise \\
\hline
\end{longtable}

\subsection{Quel outil utiliser}

\begin{itemize}
  \item Je veux recuperer un identifiant dans l'URL -> \verb|$_GET|
  \item Je veux recuperer une valeur d'un formulaire -> \verb|$_POST|
  \item Je veux garder une donnee d'une page a l'autre -> \verb|$_SESSION|
  \item Je veux envoyer l'utilisateur vers une autre page -> \verb|header('Location: ...')| suivi de \verb|exit|
  \item Je veux afficher plusieurs cartes en grille -> \verb|row row-cols-... g-4| avec des \verb|col|
  \item Je veux un champ de formulaire bien stylise -> \verb|form-control|
  \item Je veux regrouper donnees et actions -> une classe (\verb|class|)
\end{itemize}

% ============================================================
\section{Regles d'evolution de ce document}
% ============================================================

\begin{itemize}
  \item Ce document suit les memes regles que la base de reference
  Dart/Flutter : pas de retour a zero, mise a jour progressive, statut
  honnete de chaque notion.
  \item Quand une nouvelle page ou une nouvelle classe est ajoutee au projet
  web, les notions correspondantes sont ajoutees ici.
  \item Quand le projet passera a une vraie base de donnees, une nouvelle
  section sera ajoutee pour PDO ou l'outil choisi, sans supprimer ce qui est
  deja documente sur les sessions.
\end{itemize}

\end{document}
